\section*{v2.6.0 Lyapunov Ledger Closure and Forward-Invariant Regimes}
\addcontentsline{toc}{section}{v2.6.0 Lyapunov Ledger Closure and Forward-Invariant Regimes}

\begin{quote}
\textbf{Normalization note (v2.17.1 local review).}
This block is retained as a \emph{historical closure-stability and Lyapunov precursor}. Its role is supporting and developmental; the normalized proof-bearing closure chain is carried by the unfolding phases.
\end{quote}

The next logical step is to turn the budget grammar into a genuine closure device.
As long as one only writes transport inequalities, the proof still depends on a
stagewise reading. For the structure-first programme this is not yet strong enough.
One wants a single scalar functional that certifies that the admissible sector is
preserved under the whole transport chain. This is the role of the Lyapunov-ledger
closure introduced in v\docversion.

\subsection*{1. Weighted structural energy}
Choose positive weights
\[
\omega_j>0,
\qquad
\omega^{N}_j>0,
\qquad
\omega_R>0,
\qquad
\omega_\Gamma>0,
\]
and define the weighted structural energy
\[
\mathcal L
:=\sum_j \omega_j B_j + \sum_j \omega^{N}_j N_j + \omega_R\mathfrak R + \omega_\Gamma\Gamma.
\]
The aim is to choose the weights so that cross-channel transfer and residual spill
can be absorbed into one dissipative inequality. In this language, the ledger is no
longer just a bookkeeping device; it becomes the carrier of a monotonicity law.

\paragraph{Definition 1 (Lyapunov-admissible weight system).}
A family of weights is called Lyapunov-admissible if, after inserting the transport,
absorption, and source inequalities from the previous section, one obtains constants
$\kappa>0$ and $C_{\mathrm{tag}},C_{\mathrm{match}}\ge0$ such that
\[
\frac{d}{d\tau}\mathcal L
\le -\kappa \mathcal L
+ C_{\mathrm{tag}}\,\mathrm{Res}_{\mathrm{tag}}
+ C_{\mathrm{match}}\,\Delta_{\mathrm{match}}.
\]
This means that every admissible leakage path is paid for either by global decay or
by an explicitly tagged source of unresolved structure.

\subsection*{2. Weighted closure theorem}
The abstract transport statements from v2.5.0 are now compressed into a single
persistence theorem.

\paragraph{Theorem 1 (weighted ledger closure).}
Assume that the channel transport matrix has a strict damping gap, that every active
normal channel is absorption-efficient, and that residual production is source-tagged.
Then there exists a Lyapunov-admissible weight system for which
\[
\frac{d}{d\tau}\mathcal L
\le -\frac{\kappa}{2}\mathcal L
+ C_{\mathrm{tag}}\,\mathrm{Res}_{\mathrm{tag}}
+ C_{\mathrm{match}}\,\Delta_{\mathrm{match}}.
\]
In particular, if the tagged residual and probe mismatch stay integrably controlled,
then the whole derivation remains in a forward-invariant admissible regime.

\begin{proof}
Choose the weights recursively along the dominance order induced by the transfer
matrix, starting from the most dissipative channels and moving toward the weaker
ones. Because the damping gap is strict, the transfer contributions can be absorbed
by enlarging downstream weights finitely many times. The absorption-efficient normal
channels contribute additional negative terms for the dangerous normal mass. Residual
production enters only through tagged sources and therefore cannot create an
unidentified positive contribution. After summing the weighted inequalities and
absorbing the remaining off-diagonal terms, one obtains the stated estimate.
\end{proof}

This is the first point at which the document obtains a real forward-invariance
statement rather than only a local transport picture. The admissible proof sector is
now a dynamically protected cone.

\subsection*{3. Forward-invariant regime and restart stability}
One practical advantage of the Lyapunov-ledger closure is that it explains why the
proof can be resumed after interruptions without re-deriving the entire local chain.
If at some stage $\tau_0$ the state satisfies
\[
\mathcal L(\tau_0)\le M,
\qquad
\int_{\tau_0}^{\tau_1}\mathrm{Res}_{\mathrm{tag}}\,d\tau<\infty,
\qquad
\int_{\tau_0}^{\tau_1}\Delta_{\mathrm{match}}\,d\tau<\infty,
\]
then Gronwall's inequality implies a bound of the form
\[
\mathcal L(\tau)
\le e^{-\kappa(\tau-\tau_0)/2}M
+ \int_{\tau_0}^{\tau}e^{-\kappa(\tau-s)/2}
\bigl(C_{\mathrm{tag}}\,\mathrm{Res}_{\mathrm{tag}}(s)+C_{\mathrm{match}}\,\Delta_{\mathrm{match}}(s)\bigr)\,ds.
\]
Hence restarting the argument from an intermediate checkpoint is legitimate as long
as the checkpoint data are reported in ledger form. This gives a rigorous version of
proof memory: the document needs to preserve the weighted state, not the full local
micro-history.

\subsection*{4. Probe inheritance of invariance}
The two-probe picture inherits the same closure principle. Define the weighted probe
functional
\[
\mathcal L_{\mathrm{probe}}
:= \alpha_1 B^{(1)}_{\mathrm{tri}} + \alpha_2 B^{(2)}_{\mathrm{tri}}
+ \alpha_Q B^{(Q)}_{\perp} + \alpha_M B_{\mathrm{match}}
+ \alpha_R B^{\mathrm{probe}}_{\mathrm{res}}.
\]
If the probe transfer system has a strict damping gap and the mismatch source remains
ledger-tagged, then the same weighting argument yields
\[
\frac{d}{d\tau}\mathcal L_{\mathrm{probe}}
\le -\kappa_{\mathrm{probe}}\mathcal L_{\mathrm{probe}}
+ C^{\mathrm{probe}}_{\mathrm{tag}}\,\mathrm{Res}^{\mathrm{probe}}_{\mathrm{tag}}
+ C^{\mathrm{probe}}_{\mathrm{match}}\,\Delta_{\mathrm{match}}.
\]
Thus the probe sector is not merely compatible with the main ledger; it inherits the
same forward-invariant discipline. The search geometry therefore remains structurally
confined rather than opening an uncontrolled side channel.

\subsection*{5. Structural consequence for later field equations}
The future field-equation layer can now be read more sharply. A field equation will
not merely translate channel budgets into continuous variables; it must preserve the
forward-invariant regime induced by $\mathcal L$. In other words, admissible field
variables should be chosen so that the induced energy identity mirrors the weighted
ledger closure. This provides a direct criterion for accepting or rejecting later
field-equation ansaetze.

\subsection*{6. Consequence for appendix tools}
The tool appendix can be refined once more. Besides the roles introduced in v2.5.0,
a strong tool may now be marked as one of the following:
\begin{enumerate}[label=(\alph*)]
  \item weight constructor,
  \item off-diagonal absorber,
  \item invariant-region verifier, or
  \item restart certificate.
\end{enumerate}
This is useful because it identifies precisely which lemmas contribute to the global
Lyapunov closure and which only support local transport estimates.

\subsection*{7. v\docversion\ readiness criterion}
The present step should be regarded as successful if the following can now be done:
\begin{itemize}
  \item the entire admissible sector can be encoded in one weighted ledger functional;
  \item restart points can be justified by checkpoint data in ledger form;
  \item probe arguments inherit the same invariant-region control as the main proof;
  \item later field equations can be tested against a precise closure criterion; and
  \item appendix tools can be sorted by whether they contribute to local transport or global invariance.
\end{itemize}

This is the gain of v\docversion. The companion document now carries not only staged
closure, channel specialization, and budget transport, but also a genuine
forward-invariant structural regime that stabilizes the proof under continuation,
probe comparison, and later equation-level readout.


